% status: ZERO
% Printed Chapter 37 / stable source ch20. Rules: AUTHORING.md.
\chapter{Structured Output at Engine Speed}\label{ch:structured-output}

\begin{ChapterIntent}
Structured generation is not a better prompt. It is a second state machine attached
to each decoding request. At every committed token the engine must intersect the
model's distribution with the continuations accepted by a grammar, schema, regex or
tool protocol, then advance constraint state exactly once. This chapter follows that
state from schema compilation through tokenizer-aware automata, mask construction,
sampling and commit; explains why a 16-KB mask can cost more than a model forward
pass; and connects constrained decoding to batching, cancellation, prefix reuse and
the speculative-decoding contract developed later.
\end{ChapterIntent}

\OpeningSection{When checking the format costs more than running the model}

The edition's M1 measurement record asks Llama~3.2~3B for a fictional-book JSON
record twenty times. Unconstrained, 19 of 20 responses parsed and matched the desired
schema; one wrapped otherwise valid JSON in a Markdown fence. With the schema enforced
by llama.cpp, 20 of 20 were valid. Median decode cost was 40.1 ms/token constrained
versus 39.5 ms/token unconstrained---no meaningful overhead at that measurement scale.

Then one field changed. Each genre had to be one of 2,000 synthetic labels,
\texttt{genre-0000} through \texttt{genre-1999}. Validity remained 20 of 20, but
median decode cost rose to about 148 ms/token, 3.7 times the ordinary constrained run
(measured; \texttt{research/measurements/2026-09-24-m1-part3.md}). The grammar had
become more expensive than the model step it gated.

The pinned llama.cpp sampling path explains the qualitative difference. It can first
sample cheaply and test the selected token; when that candidate violates the grammar,
it falls back to applying grammar constraints across the candidate set before
resampling. With an ordinary JSON schema, the model already prefers punctuation and
field structure that are often legal. With a large synthetic enumeration, its natural
continuations disagree with the allowed language much more often.

Structured-output overhead is therefore not one constant. It depends on vocabulary
size, parser state, constraint shape, model--grammar agreement, compilation reuse and
where mask computation sits on the CPU/GPU critical path.

\Foundations{Validity is a language property, not a truth property}{
A constraint defines a language $L$ of permitted output strings. Constrained decoding
guarantees that the committed output remains on a prefix that can lead to a member of
$L$, and generation may terminate only in an accepting state. This can guarantee
syntax such as JSON shape, an enumeration or a tool-call envelope. It does not prove
that a year is historically correct, a URL is safe, a tool should be called, or an
argument is authorized. Those are semantic/application policies above the grammar.
Tokenization is part of the correctness boundary because the model chooses token IDs,
not characters. Chapter~\ref{app:tokens} explains the byte/token boundary;
Chapter~\ref{app:sampling} explains how a permitted set changes sampling.
}

\NapkinSection{Napkin math: the mask is tiny; discovering it is expensive}

Let vocabulary size be $V$, active constrained sequences $B$, and average cost to test
one token against one parser state $t_c$. A naive implementation costs
\begin{equation}
T_{\rm check}\approx B V t_c.
\label{eq:structured-naive}
\end{equation}
For $V=128{,}256$, $B=64$ and $t_c=1\,\mu$s, this is 8.21 CPU-seconds of token checks
per decode iteration (illustrative). Clearly those checks must be cached, vectorized,
parallelized, avoided or overlapped.

By contrast a dense validity bitset costs only
\[
M_{\rm mask}=\frac{V}{8}.
\]
At $V=128{,}256$, that is 16,032 bytes per sequence, or about 1.0 MiB for 64 sequences
(derived). The dominant problem is usually constructing/updating the allowed set and
moving/applying it at the right time, not storing the bits.

If a compiled state has $C$ context-dependent candidates and the other $V-C$ token
decisions are precomputed, the dynamic term becomes roughly
\[
T_{\rm dynamic}\approx B C t_c.
\]
The engineering goal is to make $C\ll V$ and hide the remaining work behind model
execution without using stale constraint state.

\section{The engine contract}

For request $r$, let $g_r$ denote constraint state and $z_r$ the model state/KV after
the committed token prefix. One decode transaction is:

\begin{enumerate}
\item run the model from committed $z_r$ and obtain logits $\ell$;
\item compute allowed token set $A(g_r)$;
\item set $\ell_i=-\infty$ for $i\notin A(g_r)$;
\item sample/select token $x$ from the remaining distribution;
\item compute candidate next grammar state $g'_r=\delta(g_r,\mathrm{bytes}(x))$;
\item commit token and $g'_r$ together, recording the valid model-state frontier;
\item permit EOS only if the grammar is in an accepting configuration.
\end{enumerate}

The ordering is a correctness contract. Sampling from the previous state's mask and
advancing the grammar afterward is not equivalent. Nor may a failed/cancelled model
step advance $g_r$ while KV remains uncommitted. A sampled token's own KV
usually does not exist yet: it is created when that token becomes an input
to the next forward pass. The transaction aligns state with the legal
history, not necessarily with identical numeric frontier counts. Track
the most recent committed token separately from the most recent position
whose target KV has actually been computed.

\EngineFigure{ch20-transaction}{Mask construction consumes the current grammar state. Only the selected committed token advances grammar state; the target-KV frontier is tracked separately until that token is fed into the model.}{fig:structured-transaction}

\section{Compile the specification before decoding}

Applications usually provide JSON Schema, a regular expression, an EBNF-like grammar,
or a tool/function schema. The runtime should not reinterpret the full high-level
document at every token. It compiles the specification into a representation suited
to incremental recognition.

Regular languages can use deterministic or nondeterministic finite automata. Nested
JSON/CFG constructs need stack-bearing machinery such as a pushdown automaton or a
parser representation with equivalent expressive power. The compiled object also
contains tokenizer-specific acceleration data: token byte strings, state-dependent
transition tables, precomputed masks or categories of tokens whose validity can be
decided without consulting deep stack context.

Compilation belongs in request admission/accounting because it consumes CPU time and
memory before the first constrained token. A service with thousands of one-off schemas
can be compilation-bound even if steady-state mask generation is fast.

\section{Bytes, Unicode and token IDs}

The grammar's alphabet and the tokenizer's emitted representation must agree. Many
tokenizers expose byte-level pieces internally; visible Unicode characters can span
multiple UTF-8 bytes, and one model token can contain several characters or punctuation
crossing grammar boundaries.

Suppose the grammar expects JSON fragment \texttt{"x":}. A vocabulary token might
encode \texttt{"x"} alone, another \texttt{":}, and another an even larger piece
including the colon. Correctness requires feeding the entire token contribution
through the recognizer from the current state. Testing only its first character admits
tokens whose suffix violates the grammar; testing only whole visible characters can
mishandle byte-fallback tokens.

The safe abstraction is
\[
\delta^*(g,b_1b_2\ldots b_n),
\]
the parser state after consuming every byte/piece contributed by token $x$. The token
is legal only if the transition exists and leaves a prefix from which acceptance is
possible.

\section{A hand-worked automaton}

Consider the tiny language
\[
L=\{\texttt{\{"ok":true\}},\ \texttt{\{"ok":false\}}\}.
\]
After the fixed prefix \texttt{\{"ok":}, the parser is at state $q_v$. From $q_v$,
the only valid textual branches are \texttt{true} and \texttt{false}. If the tokenizer
contains tokens for \texttt{true}, \texttt{false}, \texttt{tr}, \texttt{fa},
\texttt{true\}}, and \texttt{null}, then the first five can be legal depending on how
their complete byte strings advance the state; \texttt{null} cannot.

After choosing \texttt{tr}, the next state admits continuations that complete
\texttt{true}. The mask is therefore a function of parser state, not merely of the
schema. This tiny example is implemented in the chapter lab.

\section{Context-free state and the persistent stack}

JSON arrays and objects nest. A recognizer must remember which constructs remain open,
which properties are legal, and what continuation is expected after each nested
value. Copying a deep mutable stack for every candidate token or speculative branch is
wasteful.

XGrammar uses persistent stack structures so branches share unchanged history. This
supports cheap trial transitions and rollback while retaining context required for CFG
execution. Its other key optimization divides vocabulary work into
context-independent decisions that can be prechecked and a much smaller dynamic set
that needs runtime interpretation. The authors report up to 100$\times$ speedup over
their compared structured-generation engines and near-zero serving overhead in
evaluated integrations. Those are authors' measurements, not universal engine costs.
\cite{dong2025xgrammar}

\section{Mask representations}

An allowed set can be represented as:
\begin{itemize}
\item a dense bitset of $V$ bits;
\item a sparse vector of allowed token IDs;
\item a sparse vector of forbidden IDs;
\item ranges/classes plus exceptions;
\item a device-resident mask updated from compact CPU deltas.
\end{itemize}

Which wins depends on mask density and sampling backend. If only 12 tokens are legal,
shipping 12 IDs can be cheaper than writing a 25-KB mask for a 200K vocabulary. If
190K are legal, a forbidden sparse list or dense bitset is better. The representation
should be chosen from measured density and application cost, not aesthetic preference.

\section{Apply the mask without serializing the GPU}

A naive serving loop is:
\[
\text{GPU forward}\rightarrow\text{CPU mask}\rightarrow\text{copy mask}
\rightarrow\text{GPU sample}\rightarrow\text{next forward}.
\]
Every arrow can become a synchronization point. Fast implementations precompute what
they can, overlap CPU grammar work with accelerator execution, reuse device buffers,
and fuse mask application with sampling when practical.

The dependency cannot be wished away: the exact next grammar state depends on the
committed token. Pipelining must preserve that dependency. Work for request $r$ may be
prepared once $x_t$ is known while the accelerator performs independent work for other
requests or the next model invocation where ordering permits.

\section{Continuous batching means every row has different grammar state}

In a continuous batch, row 0 may be inside a JSON string, row 1 choosing a property,
row 2 unconstrained, and row 3 inside a 2,000-way enum. One global mask is incorrect.

The scheduler therefore carries per-request constraint handles alongside KV/block
tables and sampling parameters. A batch builder emits a mapping from model rows to
constraint states and mask buffers. When a request finishes, cancels or is preempted,
its grammar-state ownership follows the same lifecycle discipline as other request
state.

Constraint CPU cost can also affect scheduling. A batch containing many expensive
dynamic grammars may leave the GPU waiting even though token/KV budgets look healthy.
An industrial scheduler should observe grammar compilation and mask latency rather
than treating constrained tokens as identical to unconstrained ones.

\section{Prefix caching does not cache parser state automatically}

Two requests can share an identical prompt prefix and therefore reuse KV while asking
for different output schemas. The KV cache proves equality of model computation for
the prefix; it says nothing about output-language state.

For most API patterns the grammar begins at generation, so both requests initialize
their independent grammar state after the prompt. If constrained generation itself is
resumed from a cached generated prefix, reuse requires a compatible constraint
namespace and a parser checkpoint corresponding to that exact output prefix. Never
infer grammar-state equivalence from KV identity alone.

\section{Cancellation and preemption}

Constraint state is cheap compared with KV, but it still has ownership. Cancellation
must retire it with the request. Recomputation preemption can rebuild model KV from
tokens; grammar state can likewise be reconstructed by replaying committed output
through the compiled recognizer, or preserved as a compact checkpoint.

The invariant is:
\[
g_r=\delta^*(g_{0,r},\mathrm{committed\ output\ bytes}).
\]
A stored parser checkpoint is an optimization. The committed token sequence is the
source of truth.

\section{Jump-forward decoding}

A grammar sometimes has only one legal textual continuation for several characters.
Running a full model decode step for each forced punctuation token spends compute on a
choice that does not exist. Jump-forward techniques append a forced span and process
it efficiently, often as prefill-like work.

The optimization has two hazards. First, a forced \emph{text} span must be tokenized
compatibly with model context; an arbitrary retokenization can change the model's
token history. Second, a supposedly deterministic grammar path may still admit
multiple tokenizations even when it admits one byte string. The engine must reason
about token-level legality, not merely print the characters.

SGLang's structured-generation work uses compressed finite-state machinery and
jump-forward decoding as part of this optimization family.

\section{Sampling after masking}

Given original probabilities $p(x)$ and legal set $A$, ordinary token masking samples
\[
p_{\rm local}(x)=\frac{p(x)\mathbf{1}[x\in A]}
{\sum_{y\in A}p(y)}.
\]
This is locally normalized. It does not generally equal the distribution over complete
strings obtained by conditioning the original model on eventual grammaticality.

Grammar-Aligned Decoding formalizes this mismatch. A locally attractive legal token
may lead into a region where all eventual grammatical completions have very low model
probability; a different token with lower immediate probability may lead to much more
probable valid futures. Park et al. propose an adaptive method targeting the
conditioned distribution. \cite{park2024gad}

This is not a reason to avoid constraints. It tells us what guarantee masking provides:
validity, not distributional neutrality.

\section{Use constraints for structure, not fabricated knowledge}

The 2,000-genre experiment illustrates a product mistake. A large enum forces the
model to choose one legal symbol even when none expresses what it believes. The result
is syntactically perfect and semantically arbitrary.

Good uses include JSON punctuation, required keys, tool names from an authorized
registry, protocol envelopes and finite control codes. Riskier uses force substantive
answers into arbitrary enumerations without an ``unknown''/free-text escape. Schema
design is part of model-quality engineering.

\section{Tool calls add a security boundary}

A grammar can guarantee that
\[
\{\texttt{"tool":"transfer","amount":100}\}
\]
parses. It cannot authorize the transfer. Tool identity, argument types, permissions,
limits and user confirmation remain application/policy checks after parsing.

Structured generation narrows syntax; it must not collapse syntax validation and
capability authorization into one layer. This distinction becomes critical for
agentic systems whose structured output directly triggers side effects.

\section{Dynamic schemas change the economics}

Classic structured-generation optimization assumes a schema is reused enough to
amortize compilation. Agent systems can construct tool menus and conditional grammars
per turn. XGrammar 2 targets this workload with dynamic dispatch semantics, JIT
compilation, cross-grammar caching and broader parser machinery; its authors report
more than 6$\times$ faster grammar compilation than compared engines on
evaluated dynamic workloads, not a sixfold end-to-end token-generation speedup.
\cite{li2026xgrammar2}

The architectural lesson is that compilation cache identity now matters. A compiled
artifact should be keyed by canonical constraint semantics, tokenizer/version and
backend options, not merely by an untrusted schema string pointer. Cache memory needs
bounds and eviction just like prefix KV.

\section{Structured output meets speculative decoding}

Chapter~\ref{ch:speculative-decoding} will draft several candidate tokens before the
target model verifies them. A grammar cannot simply advance through the entire draft
and hope verification accepts it. Each drafted token must be legal along the candidate
constraint-state path, and rejection must roll grammar state back to the last committed
token.

Persistent parser state makes this natural:
\[
g_0\rightarrow g_1\rightarrow\cdots\rightarrow g_k
\]
is a candidate branch. Verification commits a prefix of these states and discards the
rest. This is the same transaction principle used for KV and chunked-prefill progress.

\section{What to measure}

For every structured request class record:
\begin{itemize}
\item schema/grammar compilation latency and cache-hit rate;
\item dynamic mask-construction latency per token;
\item allowed-set cardinality or density;
\item CPU time waiting on grammar work;
\item host-to-device mask bytes;
\item decode step time constrained versus unconstrained;
\item fallback/full-vocabulary check frequency where applicable;
\item output validity and semantic task quality separately.
\end{itemize}

Report percentiles, not only means. One pathological parser state can dominate tail
ITL. Also stratify by schema family: ``JSON schema'' covers everything from a fixed
three-field object to thousands of enums and recursive structures.

\section{Failure tests that distinguish real correctness}

A useful suite includes a token whose prefix is legal but suffix is illegal; a UTF-8
character split across byte-fallback tokens; deeply nested arrays/objects; EOS in a
nonaccepting state; a schema with an empty legal set; two requests sharing KV but not
grammar; cancellation after mask construction but before token commit; and a
speculative branch that is partly rejected.

Differentially compare the optimized mask builder with a brute-force oracle that feeds
every token's complete bytes through an independent recognizer. Performance tests do
not substitute for this oracle: a fast wrong mask silently changes the language.

\EngineFigure{ch20-batch-pipeline}{A continuous batch carries per-request constraint state. Compilation is amortized where possible; mask work is overlapped, but each row samples against the mask derived from its own committed grammar state.}{fig:structured-batch}

\EngineFigure{sys-grammar-sequence}{UML sequence view of grammar-constrained decoding (conceptual). Selection uses the intersection of the model distribution and the currently allowed set; the constraint state advances with the committed token. No particular CPU/GPU placement or parallel schedule is implied.}{fig:sys-grammar-sequence}

\LedgerSection
\begin{ConstraintLedger}
Correctness & buys & Outputs remain inside the declared syntactic language and terminate only in accepting states.\\
Compute & spends & Parser transitions and mask construction can approach $BV$ token checks without acceleration.\\
Memory bandwidth & spends & Dense masks or mask deltas must be written and possibly transferred/applied each step.\\
Latency & trades & Precomputation and overlap can make ordinary schemas nearly free; pathological dynamic constraints can dominate ITL.\\
Capacity & spends & Compiled grammars, persistent stacks and cached masks consume bounded CPU/device memory.\\
Correctness & risks & Token/byte mismatch, stale parser state, premature EOS, rollback bugs and local distribution distortion.\\
Cost & trades & Valid-by-construction output removes retry/parsing waste; expensive one-off schemas can consume more CPU than expected.\\
\end{ConstraintLedger}

\LabSection{token-aware masking with an independent oracle}

The Rust lab in \texttt{code/labs/ch20-structured-output/} deliberately uses a tiny
vocabulary so every transition can be inspected. It recognizes exactly
\texttt{\{"ok":true\}} and \texttt{\{"ok":false\}}. Vocabulary pieces include complete
values, partial values, punctuation-crossing pieces and an invalid \texttt{null}.

The optimized mask builder memoizes allowed token IDs per parser state. The independent
oracle does not consult that cache or its prefix-state transition table: it
concatenates each token's bytes to the committed bytes and directly checks
whether either of the two complete permitted strings starts with that result.
Tests compare both masks at
every reachable state, reject premature EOS, and verify rollback by deriving state
again from committed bytes.

The lab is intentionally not a benchmark of XGrammar, llama.cpp or GPU serving. It
teaches the state/commit contract that an optimized backend must preserve.

\HappenedSection{grammar libraries moved inside the engine}

Structured output evolved from application-side retry loops into a serving primitive.
llama.cpp has long supported grammar-constrained sampling; the pinned measurement
shows both its cheap case and a large-enum slow path. SGLang developed structured
generation and jump-forward techniques. XGrammar co-designed grammar execution with
LLM serving using context-independent token prechecks, persistent stack state and
overlap with GPU execution. Modern engines integrate specialized structured-output
backends rather than interpreting JSON Schema from scratch after every token.

The resulting architecture resembles the rest of this book: compile expensive static
work once, keep a compact per-request state machine, batch the dynamic part, avoid
synchronization, and make rollback/ownership explicit.

\FrontierSection{2026-10-05}

XGrammar remains an important reference for high-performance CFG execution; its paper
reports up to 100$\times$ acceleration over compared structured-generation systems and
near-zero overhead in integrated serving. Grammar-Aligned Decoding demonstrates that
ordinary local masking can distort the model's distribution over valid complete
outputs. JSONSchemaBench broadens evaluation from toy schemas to a 10K-schema corpus
and emphasizes efficiency, coverage and output quality as separate axes.

The 2026 frontier is dynamic agentic structure. XGrammar 2 introduces dynamic dispatch,
JIT compilation and cross-grammar caching for tool-rich workloads, reporting more than
6$\times$ faster grammar compilation than compared structured-generation
engines in its evaluation.
These are research-system measurements; production claims require pinned engine source
and deployment measurements.

\ExercisesSection
\begin{enumerate}
\item \emph{Derivation.} With $V=200{,}000$, $B=64$ and an 8-ms model step, what
average $t_c$ would keep naive equation~\eqref{eq:structured-naive} below 10\% of the
step? Explain why the answer motivates precomputation.
\item \emph{Trace.} For the lab grammar, trace parser states after tokens
\texttt{\{"ok":}, \texttt{tr}, \texttt{ue\}} and identify where EOS first becomes
legal.
\item \emph{Representation.} Compare bytes written by a dense bitset with a sparse
list when 12 of 200,000 token IDs are legal. State the ID width you assume.
\item \emph{Implementation.} Add speculative branches to the lab. Acceptance:
candidate grammar states are persistent/cloneable; rejecting token $j$ restores the
state after committed token $j-1$; oracle parity holds after every rollback.
\item \emph{Security.} Give an example of a syntactically valid tool call that must
still be rejected by authorization policy.
\item \emph{Falsification.} Construct a grammar where locally normalized masking
chooses a different highest-probability path than ranking complete valid strings by
the original model probability.
\end{enumerate}

\Needspace{18\baselineskip}
\section{Worked problems}

\Needspace{12\baselineskip}
\subsection{Naive check budget}
\textbf{Problem.} \emph{Derivation.} With $V=200{,}000$, $B=64$ and an 8-ms model step, what
average $t_c$ would keep naive equation~\eqref{eq:structured-naive} below 10\% of the
step? Explain why the answer motivates precomputation.

\textbf{Solution.} Ten percent of 8 ms is 0.8 ms. Equation~\eqref{eq:structured-naive} gives
$t_c\le0.0008/(64\cdot200000)=6.25\times10^{-11}$ s, or 0.0625 ns/check.
That is not a realistic scalar parser budget; the vocabulary work must be avoided,
amortized, parallelized or overlapped.

\Needspace{12\baselineskip}
\subsection{Sparse versus dense}
\textbf{Problem.} \emph{Representation.} Compare bytes written by a dense bitset with a sparse
list when 12 of 200,000 token IDs are legal. State the ID width you assume.

\textbf{Solution.} A 200,000-bit dense mask is 25,000 bytes. Twelve legal 32-bit token IDs need only
48 bytes before container metadata, so sparse allowed IDs win strongly at that state.

\Needspace{12\baselineskip}
\subsection{Tool-call security}
\textbf{Problem.} \emph{Security.} Give an example of a syntactically valid tool call that must
still be rejected by authorization policy.

\textbf{Solution.} A grammar may accept a transfer tool with a numeric amount while authorization rejects
it because the caller lacks permission, the amount exceeds a limit, or confirmation is
missing. Syntax and authority are independent guarantees.

\Needspace{12\baselineskip}
\subsection{EOS}
\textbf{Problem.} \emph{Trace.} For the lab grammar, trace parser states after tokens
\texttt{\{"ok":}, \texttt{tr}, \texttt{ue\}} and identify where EOS first becomes
legal.

\textbf{Solution.} EOS is legal only after the recognizer reaches an accepting state representing the
complete closing brace; accepting a prefix such as \texttt{\{"ok":tr} would violate
the language even though it can still be extended to a valid string.

